\section{General conventions and notation}\label{sec:conventions}


We collect the conventions used throughout. This is a reference
section for the written notes, not a preliminary lecture. Each main
lecture repeats the minimum notation and equations it needs. The
conventions follow \cite{McAllister:2025qwq}, \S\,2.2, and have
been streamlined for the present audience.

\subsection{Indices and spacetime}

The ten dimensions of the ambient string-theoretic spacetime are
indexed by capital Latin letters $M,N,\ldots = 0,\ldots,9$. The
ansatz of physical interest splits these into four ``external''
dimensions $\mu,\nu,\ldots = 0,1,2,3$ corresponding to the
four-dimensional spacetime that our cosmology lives in, and six
``internal'' dimensions $m,n,\ldots = 4,\ldots,9$ corresponding to a
compact real $6$-manifold $X$:
\begin{equation}\label{eq:ansatz}
    \mathrm{d}s^2_{10} \;=\; \mathrm{e}^{2A(y)}\, g^{(4)}_{\mu\nu}(x)\,\mathrm{d}x^\mu \mathrm{d}x^\nu
        \,+\, \mathrm{e}^{-2A(y)}\,g^{(6)}_{mn}(y)\,\mathrm{d}y^m \mathrm{d}y^n \,.
\end{equation}
The factor $\mathrm{e}^{2A(y)}$ is the \emph{warp factor}; we will set it to
unity at first and reinstate it in Lecture~3. We use signature
$(-,+,+,\ldots,+)$ for $g^{(10)}$.

\subsection{Differential forms and the Hodge decomposition}

For a real $C^\infty$ manifold $M$ of dimension $D$ we write
$\Omega^r(M)$ for the space of smooth $r$-forms,
$\mathrm{d}:\Omega^r\to
\Omega^{r+1}$ for the exterior derivative, and $\star:\Omega^r\to
\Omega^{D-r}$ for the Hodge star with respect to a Riemannian metric
$g$. The adjoint
$\mathrm{d}^\dagger = (-1)^{Dr+D+1}\star \mathrm{d}\,\star$ and the
Laplacian
$\Delta = \mathrm{d}\mathrm{d}^\dagger
+ \mathrm{d}^\dagger\mathrm{d}$ act on $\Omega^r(M)$.

We denote by $\mathcal{H}^r(M)$ the space of harmonic $r$-forms. For
$M$ compact orientable, Hodge's theorem identifies
\begin{equation}\label{eq:hodge-iso}
    H^r(M,\mathbb{R}) \;\cong\; \mathcal{H}^r(M) \,.
\end{equation}
We write $b_r(M) = \dim_\mathbb{R} H^r(M,\mathbb{R})$ for the Betti
numbers.

\subsection{Complex manifolds and Dolbeault cohomology}

A \emph{complex manifold} of complex dimension $k$ is a real
$2k$-manifold with complex coordinate charts whose transition
functions are holomorphic. The tangent bundle splits
$TM\otimes\mathbb{C} = T^{1,0}M\oplus T^{0,1}M$, and differential forms
decompose into $(p,q)$-types,
\begin{equation}
    \Omega^r(M,\mathbb{C}) \;=\;
        \bigoplus_{p+q=r}\,\Omega^{p,q}(M) \,,
\end{equation}
with $d = \partial + \bar\partial$ on a Hermitian manifold and
$\bar\partial^2 = 0$. The $\bar\partial$-Laplacian
$\Delta_{\bar\partial} = \bar\partial\bar\partial^\dagger
+ \bar\partial^\dagger\bar\partial$ defines harmonic
$(p,q)$-forms; we write $\mathcal{H}^{p,q}_{\bar\partial}(M)$ for the
space and
$h^{p,q}(M) = \dim \mathcal{H}^{p,q}_{\bar\partial}(M)$
for the \emph{Hodge numbers}.

On a K\"ahler manifold (Hermitian metric whose K\"ahler form $J$ is
closed) one has the identity
$\Delta_d = 2\Delta_{\partial} = 2\Delta_{\bar\partial}$, hence
\begin{equation}\label{eq:betti-hodge}
    b_r(M) \;=\; \sum_{p+q=r} h^{p,q}(M) \,,
    \qquad h^{p,q} \;=\; h^{q,p} \,,
    \qquad h^{p,q} \;=\; h^{k-p,k-q} \,.
\end{equation}

\subsection{Calabi--Yau threefolds}

There is more than one convention for ``Calabi--Yau''; we record
ours explicitly. \emph{In these notes a Calabi--Yau $n$-fold is a
compact K\"ahler manifold $X$ of complex dimension $n$ with
trivial canonical bundle $K_X \cong \mathcal{O}_X$ and
$\pi_1(X) = 0$.} The convention matches that used in the
candidate-dS literature
\cite{McAllister:2025qwq,McAllister:2024lnt}; the
simple-connectedness clause is convenient but not standard across
the differential-geometry literature (Joyce, for instance, defines a
CY $n$-fold as a compact K\"ahler manifold with $\Hol(X) =
\mathrm{SU}(n)$, which implies $\pi_1(X)$ finite but not necessarily
trivial). Under our convention, the following standard structural
facts hold; in each case the implication should be read in the
indicated direction.
\begin{itemize}[nosep,leftmargin=*]
    \item Triviality of $K_X$ is equivalent to the existence of a
    nowhere-vanishing holomorphic $(n,0)$-form
    $\Omega \in H^0(X, K_X) = H^{n,0}(X)$.
    \item Yau's theorem (Thm.~\ref{thm:yau} below) gives the
    existence and uniqueness of a Ricci-flat K\"ahler metric in
    each K\"ahler class.
    \item For this Ricci-flat metric, the restricted holonomy
    (which, under $\pi_1(X)=0$, equals the full holonomy) satisfies
    $\Hol(X) \subseteq \mathrm{SU}(n)$; equality $\Hol(X) =
    \mathrm{SU}(n)$ requires further input, namely that $X$ is
    irreducible (not a Riemannian product) and not hyperk\"ahler,
    in which case Berger's classification
    (Theorem~\ref{thm:berger}) gives $\Hol(X) = \mathrm{SU}(n)$.
\end{itemize}
For $n=3$ the Hodge diamond and the resulting Euler characteristic $\chi(X) = 2\bigl(h^{1,1}(X) - h^{2,1}(X)\bigr)$ are stated in \S\ref{ssec:L1-CY} (Proposition `CY3 Hodge diamond'); we do not repeat them here.

\subsection{Toric notation (reminder; not derived)}

We use the standard toric-geometry conventions of
\cite{cox-little-schenck:toric}. For a complete simplicial fan
$\Sigma \subset N_{\mathbb{R}}$ (with $N \cong \mathbb{Z}^d$) we write
$V_\Sigma$ for the associated toric variety. Prime toric divisors
$D_\rho$ correspond to rays $\rho\in\Sigma(1)$. The Chow group
$A^1(V_\Sigma) \otimes \mathbb{R}$ has a distinguished cone
$\KC(V_\Sigma)$ (the K\"ahler cone) dual to the Mori cone
$\MC(V_\Sigma)\coloneqq  \overline{\mathrm{NE}}(V_\Sigma)$ (the closure of
the cone generated by classes of effective $1$-cycles; see
Definition~\ref{def:moricone}). The notation $\MC$ is used in the
lectures to avoid confusing the Mori cone with the nef cone; the
standard algebraic-geometry notation is displayed at the point of
definition. For a 4-dimensional reflexive polytope
$\Delta^\circ \subset N_\mathbb{R}$ (with polar dual
$\Delta\subset M_\mathbb{R}$, $M=N^\vee$), the Batyrev construction
\cite{Batyrev:1993oya} produces a Calabi--Yau threefold $X$ as the
generic anticanonical hypersurface in the ambient toric fourfold
$V_{\Sigma_{\mathcal{T}}}$ associated to a fine, regular, star
triangulation (FRST) $\mathcal{T}$ of $\Delta^\circ$
(Appendix~\ref{app:toric-FRST}). The ambient
$V_{\Sigma_{\mathcal{T}}}$ has at worst $\mathbb{Q}$-factorial
terminal Gorenstein singularities, located in cones of codimension
$\geq 4$ in $V_{\Sigma_{\mathcal{T}}}$ (i.e.\ at isolated points),
which generically do not meet the smooth hypersurface $X$.

We write $\kappa_{abc} = D_a\cdot D_b\cdot D_c$ for the triple
intersection numbers of (pullbacks of) prime toric divisors on $X$.
For the Calabi--Yau hypersurface $X$, the K\"ahler cone is denoted
$\KC(X) \subset H^{1,1}(X,\mathbb{R})$; the \emph{extended K\"ahler
cone} $\eKC(X)$ is the union of $\KC(X')$ over all birational models
$X'$ reached by chains of flops, glued along their common flop walls
(see Proposition~\ref{prop:flop-formulas} and
Appendix~\ref{app:kcone-extended}).

The mirror Calabi--Yau threefold of $X$ is denoted $\widetilde X$
(Batyrev construction, Appendix~\ref{app:hodge-batyrev}); we do not
use the $X^\circ$ notation. Note that the tilde has two unrelated
uses in these notes: the mirror manifold $\widetilde X$ and the
extended K\"ahler cone $\eKC = \widetilde{\mathcal{K}}$.

\subsection{Conventions for field-theory quantities}

We work in $\hbar = c = 1$ units, and adopt $\Pl = 1$, where the
four-dimensional reduced Planck mass has dimensionful value
$\Pl \simeq 2.4\times 10^{18}~\mathrm{GeV}$.
A \emph{field} on a spacetime manifold is a geometric object allowed
to vary from point to point: for example a function, a differential
form, a connection, or a metric. An \emph{action} is a functional of
these fields,
\begin{equation}
    S[\Phi,g] = \int_M \mathrm{d}^{\dim M}x\,
        \sqrt{|g|}\,\mathcal{L}(\Phi,g,\nabla\Phi,\ldots),
\end{equation}
where $\Phi$ denotes all non-metric fields and $\mathcal{L}$ is the
Lagrangian density. Varying $S$ with respect to $\Phi$ gives the
field equations. These field equations are what physicists call the
\emph{equations of motion}: the differential equations satisfied by
classical field configurations that are stationary points of the
action. Varying $S$ with respect to the metric gives the stress-energy
tensor
\begin{equation}\label{eq:stress-tensor-def}
    T_{MN}
    \coloneqq  
    -\frac{2}{\sqrt{-g}}\,
    \frac{\delta S_{\rm matter+source}}{\delta g^{MN}} \, .
\end{equation}
Thus a schematic ten-dimensional action of the form
\begin{equation}
    S^{(10)}
    =
    \frac{1}{2\kappa_{10}^2}
    \int \mathrm{d}^{10}x\,\sqrt{-g}\,
        \bigl(R^{(10)}+\mathcal{L}_{\rm matter}
        +\mathcal{L}_{\rm source}\bigr)
\end{equation}
gives the Einstein equation
\begin{equation}\label{eq:einstein-stress-convention}
    R^{(10)}_{MN} - \frac12 g_{MN}R^{(10)}
        = \kappa_{10}^2\,T_{MN}.
\end{equation}
Taking the trace in $10$ dimensions rewrites this as
$R^{(10)}_{MN}=\kappa_{10}^2(T_{MN}-\tfrac18 g_{MN}T)$, with
$T=T^K{}_K$.

A four-dimensional field is called \emph{light} relative to an
effective description with cutoff $\Lambda$ if its mass is well below
$\Lambda$, so it is retained as a dynamical variable. Fields with
masses comparable to or above $\Lambda$ are integrated out and affect
the light theory only through corrections to its couplings. In a
compactification, harmonic zero modes are massless at leading order
and hence light before fluxes, branes, or quantum corrections generate
a potential.

The
ten-dimensional gravitational coupling $\kappa_{10}^2$ is normalised
so that the leading-order action takes the form
\begin{equation}\label{eq:S10-norm}
    S^{(10)} =
    \frac{1}{2\kappa_{10}^2}\!
    \int \mathrm{d}^{10}x\,\sqrt{-g_{10}}\,
        \Bigl(R^{(10)} - \tfrac{1}{2}|\partial\tau|^2_\tau
            - \tfrac{1}{2}|G_3|^2 - \tfrac{1}{4}|\widetilde F_5|^2
            + \ldots\Bigr) \,,
\end{equation}
with $\tau = C_0 + \mathrm{i} \mathrm{e}^{-\phi}$ the axio-dilaton,
$G_3 = F_3 - \tau H_3$ the complexified three-form, and
$\widetilde F_5$ the
self-dual five-form (precise form: \eqref{eq:IIB-action}). The
norm $|\partial\tau|^2_\tau$ uses the hyperbolic metric on the upper
half-plane.

\begin{remark}[A note on physics terms]
Throughout these lectures, ``supergravity'' (SUGRA) denotes the
low-energy field-theory approximation to a chosen string background:
an action-based field theory on a Lorentzian manifold of fixed
dimension, retaining the massless fields and their leading
interactions. ``Supersymmetry'' (SUSY) denotes a
$\mathbb{Z}_2$-graded extension of the Poincar\'e algebra by
spinorial generators; an $\mathcal{N}=1$ theory in $D$ dimensions
has one such spinor generator (so $4$ real supercharges in $D=4$).
``Flux'' denotes a properly quantised closed differential form, or
more precisely its integral cohomology class. In the supersymmetric
geometric approximation, a ``brane'' is represented by a calibrated
submanifold of $X$ carrying gauge / matter degrees of freedom.
Precise definitions appear in Lectures~2 and~3.
\end{remark}



